\section{第2章の解答}\label{解答:第2章}
\paragraph{2.1}
(1) 分子・分母を $n^2$ で割ると,
\[
 \frac{2n^2+3}{5n^2-n}
 =\cfrac{2+\cfrac3{n^2}}{5-\cfrac1n}
 \longrightarrow\frac25.
\]
分母の極限は $5\ne0$ なので,商の極限法則を適用できる.

(2) 根号を含む差を有理化する.
\[
 \sqrt{n^2+2n}-n
 =\frac{2n}{\sqrt{n^2+2n}+n}
 =\cfrac2{\sqrt{1+\cfrac2n}+1}
 \longrightarrow1.
\]
二つの項がそれぞれ大きくなることだけでは差の極限は分からないが,この形なら分母が $2$ に近づくことが分かる.

漸化式では,まず $1\le a_n<2$ を帰納的に示す.\,$a_1=1$ はこの条件を満たす.\,$1\le a_n<2$ なら,
\[
 1<\sqrt3\le a_{n+1}=\sqrt{2+a_n}<2
\]
だから,次の項も条件を満たす.さらに,
\[
 a_{n+1}^2-a_n^2
 =2+a_n-a_n^2=(2-a_n)(a_n+1)>0.
\]
両項が正なので $a_{n+1}>a_n$ である.したがって数列は単調増加かつ上に有界であり,有限の極限 $L$ をもつ.

$1\le L\le2$ であり,平方根の連続性から $L=\sqrt{2+L}$ となる.これを平方すると
\[
 L^2-L-2=(L-2)(L+1)=0.
\]
範囲に合うのは $L=2$ である.漸化式に極限を代入する前に,収束することを確かめた点が大切である.

\paragraph{2.2}
公比 $r=-\frac12$ は $|r|<1$ を満たすので,無限和は
\[
 S=\cfrac1{1+\cfrac12}=\frac23.
\]
最初の $n$ 項の和を $S_n$ とすると,
\[
 S-S_n=\cfrac{\left(-\cfrac12\right)^n}{1+\cfrac12},\qquad
 |S-S_n|=\frac{2}{3\cdot2^n}.
\]
この値は $n$ とともに減少する.\,$n=9$ では $\frac1{768}>\frac1{1000}$,\,$n=10$ では $\frac1{1536}<\frac1{1000}$ となる.したがって,条件を満たす最小の自然数は $n=10$ である.

\paragraph{2.3}
$x\ne0$ で
\[
 \frac{\sin2x}{x}=2\frac{\sin2x}{2x}
\]
だから,基本極限により $\lim_{x\to0}f(x)=2$ となる.連続となるにはこの極限が $f(0)=c$ と一致する必要があり,答えは $c=2$ である.

$g$ については $g(1)=0$ であり,\,$g(1+h)=|h|\to0$ なので $1$ で連続である.しかし,\,$h\ne0$ に対して差商は
\[
 \frac{g(1+h)-g(1)}h=\frac{|h|}{h}
 =\begin{cases}1&(h>0),\\-1&(h<0)\end{cases}
\]
となる.右極限と左極限が異なるため,\,$g'(1)$ は存在しない.連続性は値のつながりを調べ,微分可能性は差商の極限まで調べるので,両者は同じ条件ではない.

\paragraph{2.4}
(1) 外側の対数の微分に,内側の $1+x^2$ の微分を掛ける.
\[
 \frac{\dd}{\dd x}\log(1+x^2)=\frac{2x}{1+x^2}.
\]
$1+x^2>0$ なので,すべての実数 $x$ で定義される.

(2) 積の微分と合成関数の微分から,
\begin{align*}
 \frac{\dd}{\dd x}\{(x^2+1)\e^{-x}\}
 &=2x\e^{-x}-(x^2+1)\e^{-x}\\
 &=-(x-1)^2\e^{-x}.
\end{align*}
$x=0$ で関数値は $1$,微分係数は $-1$ だから,接線は $y-1=-x$,すなわち $y=1-x$ である.

(3) $y=x^{\sin x}>0$ とおき,対数をとると $\log y=\sin x\log x$ となる.両辺を微分して
\[
 \frac{y'}y=\cos x\log x+\frac{\sin x}{x}
\]
を得るので,
\[
 y'=x^{\sin x}\left(\cos x\log x+\frac{\sin x}{x}\right)
 \qquad(x>0)
\]
である.指数が変数を含むため,定数指数のべきの公式だけでは微分できない.

